% =============================================================================
% 1.1 Introduction
% =============================================================================
\section{Introduction}
\label{sec:ch1-intro}

Le secteur de la construction et des travaux publics constitue l'un des piliers
de l'économie, caractérisé par une compétition intense où la réponse aux appels
d'offres représente un enjeu stratégique majeur. Chaque année, des milliers
d'entreprises consacrent des ressources considérables à la préparation de
dossiers de candidature, processus qui mobilise des compétences juridiques,
techniques, financières et administratives.

L'émergence des technologies d'intelligence artificielle, et plus
particulièrement des grands modèles de langage, ouvre des perspectives inédites
pour automatiser et optimiser ce processus. Cependant, la spécificité des appels
d'offres — leur caractère hautement structuré, leur exigence de traçabilité et
de conformité, et leur dimension juridique contraignante — impose une approche
réfléchie et méthodique.

Ce chapitre dresse un panorama des connaissances actuelles dans les domaines
connexes à notre problématique. Nous procédons en trois mouvements. Le premier
caractérise le domaine d'application~: le processus des appels d'offres, sa
production documentaire et ses enjeux de capitalisation
(sections~\ref{sec:ch1-contexte} et~\ref{sec:ch1-ged}). Le deuxième parcourt
les disciplines qui outillent une réponse technique~: traitement automatique du
langage, recherche d'information, génération augmentée par récupération, bases
vectorielles, compréhension de documents structurés et résolution d'entités
(sections~\ref{sec:ch1-tal} à~\ref{sec:ch1-entites}). Le troisième discute les
architectures de systèmes et confronte notre positionnement aux solutions
existantes (sections~\ref{sec:ch1-workflows} et~\ref{sec:ch1-solutions}).

\begin{btpnote}[Périmètre de ce chapitre]
Nous distinguons soigneusement les \btpterm{concepts} des \btpterm{outils}. Un
état de l'art rend compte des modèles, des méthodes et de leurs propriétés
démontrées~; le choix d'une bibliothèque ou d'un produit relève de la
méthodologie et de l'implémentation, traitées aux chapitres~\ref{chap:methodologie}
et~\ref{chap:resultats}. Les outils cités ici le sont à titre d'illustration
d'une approche, jamais comme un choix de conception.
\end{btpnote}
